Вот чистый русский текст (без LaTeX-команд), который написан в подразделе:Теоретический статус и фиксация параметра ( P )На ранних этапах разработки модели Decoupled Observer Patch Holography (dOPH) параметр ( P ) вводился как свободный феноменологический коэффициент в проекционном операторе Π(ρ)\Pi(\rho)\Pi(\rho)
. Этот параметр определяет силу гравитационного усиления, испытываемого барионной материей вследствие механизма декуплинга наблюдателя.При систематической подгонке ранних версий модели к данным о кривых вращения галактик SPARC (175 галактик) параметр ( P ) демонстрировал высокую устойчивость и стабильно сходился к узкому интервалу от 1.630 до 1.634, при этом наилучшее подобранное значение обычно составляло P≈1.63094P \approx 1.63094P \approx 1.63094
.Дальнейший анализ показал, что это эмпирически предпочтительное значение очень близко к точному математическому выражению 8/3\sqrt{8/3}\sqrt{8/3}
. В частности,83≈1.63299,\sqrt{\frac{8}{3}} \approx 1.63299,\sqrt{\frac{8}{3}} \approx 1.63299,
что отличается от подобранного значения всего на 0.12–0.18 %.Начиная с версии Baseline v2, мы переводим параметр ( P ) из разряда свободного подгоночного коэффициента в статус фундаментальной константы теории:P≡83≈1.63299.P \equiv \sqrt{\frac{8}{3}} \approx 1.63299.P \equiv \sqrt{\frac{8}{3}} \approx 1.63299.
Мотивация из структуры теории dOPHТакой выбор не является произвольным. Значение P=8/3P = \sqrt{8/3}P = \sqrt{8/3}
 получает естественное теоретическое обоснование внутри теории Decoupled Observer Patch Holography.В основе теории лежит концепция голографического экрана (holographic screen), на котором закодирована гравитационная информация, проецируемая на трёхмерное пространство, наблюдаемое локальным наблюдателем. При введении четвёртой аксиомы — аксиомы декуплинга — добавляется локальный механизм исправления ошибок на пересекающихся областях соседних патчей без нарушения глобальной согласованности.В результате применения этого механизма пересекающиеся области патчей эффективно «вычитаются» или проецируются из домена наблюдателя. Геометрический и информационный анализ получающегося эффективного объёма патча относительно площади голографического экрана приводит к характерному отношению 8/3. Параметр ( P ) естественно возникает как квадратный корень из этого отношения.Этот геометрический фактор проявляется при последовательном учёте как голографической проекции, так и процедуры локального декуплинга. В пределе высокого фазового консенсуса (при корректном разрешении пересечений патчей) возникающее гравитационное ускорение барионной материи приобретает усиление, модулируемое именно этим фактором.Численная близость между ранее подобранным значением (P≈1.63094P \approx 1.63094P \approx 1.63094
) и теоретически мотивированным значением 8/3\sqrt{8/3}\sqrt{8/3}
 подтверждает внутреннюю согласованность модели. Разница лежит в пределах суммарной погрешности, включающей статистическую неопределённость Monte-Carlo симуляций (до 20 000 синтетических галактик) и разброс наблюдательных данных SPARC.Последствия фиксации параметраФиксация P≡8/3P \equiv \sqrt{8/3}P \equiv \sqrt{8/3}
 позволяет добиться нескольких важных улучшений:Модель переходит от чисто феноменологического описания к теории с большей предсказательной силой за счёт устранения одного свободного параметра.
Проекционный оператор Π(ρ)\Pi(\rho)\Pi(\rho)
 и сила декуплинга Fdecouple\mathbf{F}_{\rm decouple}\mathbf{F}_{\rm decouple}
 теперь строятся на основе фундаментальной константы, вытекающей из аксиом и геометрии теории.
Стабильность модели как на синтетических данных (98.8 % на 20 000 галактик), так и на реальных данных SPARC практически не изменяется, что подтверждает совместимость теоретического значения с наблюдениями.

Важно отметить, что принятие P=8/3P = \sqrt{8/3}P = \sqrt{8/3}
 отражает наше текущее наилучшее теоретическое понимание в рамках модели dOPH. Мы не утверждаем, что это абсолютная универсальная константа, независимая от модели. Это значение является следствием голографической структуры теории в сочетании с аксиомой декуплинга в её нынешней формулировке. В случае, если будущие высокоточные данные или теоретические разработки покажут статистически значимое предпочтение другого значения, модель может быть соответствующим образом уточнена.Данная фиксация является важной вехой в развитии dOPH: переходом от подгонки параметров к более аксиоматичной и геометрически обоснованной формулировке теории.Это полный текст, который соответствует английской версии.  Хочешь, чтобы я сделал его короче, строже или добавил/убрал какие-то моменты?

